\section{CommonPile：只用 permissively licensed data 可行吗}

前面的数据集大多在质量、规模和法律风险之间折中，本节把约束推到更严格的版本：只收集 permissively licensed data。CommonPile/Comma 的价值在于验证这种路线可以训练出有竞争力模型，同时也暴露可用 token 规模、license laundering 和文档级授权继承等现实边界。

\begin{figure}[H]
\centering
\includegraphics[width=0.86\textwidth]{commonpile.png}
\caption{CommonPile：收集 8TB permissively licensed data。}
\figsource{CS336 Lecture 13 官方源引用图。}
\end{figure}

\begin{figure}[H]
\centering
\includegraphics[width=0.86\textwidth]{comma-results.png}
\caption{CommonPile/Comma results：只用 permissive data 可以做到不错，但 token 规模仍是瓶颈。}
\figsource{CS336 Lecture 13 官方源引用图。}
\end{figure}

\begin{warningbox}{读图：permissive data 也有细节}
License laundering 可能把 copyrighted work 重新包装成 permissive license；collection license 不一定延伸到单个文档；合成数据若来自未授权数据训练的模型，法律状态也不清楚。只说“permissive”不够，仍需 provenance 审计。
\end{warningbox}

